% 问题二方法：共表达图、关联选择与证据核验。
\documentclass[tikz,border=0pt]{standalone}
\usepackage[UTF8,fontset=windows]{ctex}
\usepackage{amsmath}
\setmainfont{Times New Roman}
% 方法流程图的共用样式：黑白矩形框、直角箭头、宋体正文、侧边竖排分区标题。
% 坐标单位为毫米，y 轴向下；版心宽 158 mm，侧边标题占 0--6 mm，正文框占 9--158 mm。
\usetikzlibrary{arrows.meta,calc}
% ctex 默认行距系数为 1.3，图内统一取 1。
\newcommand{\fcsize}[2]{\linespread{1}\fontsize{#1}{#2}\selectfont}
\tikzset{
  fc/.style={x=1mm,y=-1mm,>=Latex,line width=0.45pt},
  box/.style={draw,rectangle,anchor=north west,inner sep=1.5mm,outer sep=0pt,align=left,
    font=\songti\fcsize{8.6}{11}},
  % 固定正文高度并顶端对齐，使同一行各框的标题对齐；框外高为 h + 3 mm。
  h/.style={execute at begin node={\begin{minipage}[t][#1][t]{\linewidth}\raggedright},
    execute at end node={\end{minipage}}},
  io/.style={box,dashed},
  key/.style={box,line width=0.9pt},
  edge/.style={->,line width=0.45pt},
  note/.style={font=\songti\fcsize{8}{9.5},fill=white,inner sep=0.6mm},
}
% 框标题
\newcommand{\bt}[1]{{\heiti #1}\\}
% 侧边分区：\sidelabel{上沿 y}{下沿 y}{竖排文字}
\newcommand{\sidelabel}[3]{%
  \draw (0.3,#1) rectangle (6.1,#2);
  \node[align=center,font=\heiti\fcsize{9.5}{11.5},inner sep=0pt]
    at (3.2,{(#1+#2)/2}) {#3};}
% 总线连接：上方节点列表、下方节点列表、总线 y、总线左端 x、右端 x
\newcommand{\bus}[5]{%
  \foreach \u in {#1} {\draw (\u.south) -- (\u.south |- 0,#3);}
  \draw (#4,#3) -- (#5,#3);
  \foreach \l in {#2} {\draw[edge] (\l.north |- 0,#3) -- (\l.north);}}

\begin{document}
\begin{tikzpicture}[fc]
\path[use as bounding box] (0,0) rectangle (158,165);
\sidelabel{0}{21}{输\\入\\与\\定\\义}
\node[io,text width=69.5mm,h=18mm] (chem) at (9,0) {\bt{化学输入}%
250款啤酒；附件1的231列中排除5个综合指标和1个实验内标，留下225个单体化合物\\
附件2提供231列之间的Spearman秩相关矩阵};
\node[io,text width=69.5mm,h=18mm] (sens) at (85.5,0) {\bt{感官目标与“关键”的含义}%
逐款原始品评小组均分转为秩；49个可归因维度\\
“关键”指控制已选化合物后的较强关联，不等于因果作用};
\sidelabel{26}{73}{筛\\选\\模\\型}
\node[box,text width=146mm,h=13mm] (graph) at (9,26) {\bt{构建化合物共表达图与模块}%
用相关系数绝对值作边权；在$1-|\rho|$距离上层次聚类。阈值$|\rho|=0.6$划分模块，0.5与0.7作敏感性检查。};
\bus{chem,sens}{graph}{23}{45.25}{121.75}
\node[box,text width=44mm,h=23mm] (sim) at (9,47) {\bt{选择规则的独立依据}%
只用真实化学矩阵植入已知的3个化合物，比较能否找回\\
避免以感官结果反选规则};
\node[key,text width=44mm,h=23mm] (select) at (60,47) {\bt{图约束前向选择}%
第一位取绝对秩相关最大者\\
第二、三位取控制已选集合后的绝对秩偏相关最大者\\
同一模块不重复入选};
\node[box,text width=44mm,h=23mm] (sig) at (111,47) {\bt{统计证据}%
置换检验评估最大关联；后两位使用保留已选变量作用的置换\\
日期组重抽的关联区间须排除零};
\bus{graph}{sim,select,sig}{44}{32.5}{134.5}
\sidelabel{80}{110}{稳\\定\\性}
\node[box,text width=69.5mm,h=24mm] (boot) at (9,80) {\bt{重复抽样}%
25个品评日期组有放回重抽1000次；每次重建图和模块，再选前三个\\
记录逐化合物和逐模块的入选频率};
\node[box,text width=69.5mm,h=24mm] (compare) at (85.5,80) {\bt{替代方法对照}%
无图约束、单次边际选择、弹性网络和梯度提升树置换重要度\\
用Jaccard交并比衡量前三个集合的一致性};
\bus{sim,select,sig}{boot,compare}{77}{32.5}{134.5}
\sidelabel{114}{165}{解\\释\\与\\交\\付}
\node[box,text width=44mm,h=25mm] (grade) at (9,114) {\bt{维度证据分级}%
显著且第一位入选频率至少0.6：较强证据\\
仅显著：有限证据\\
其余：未找到可靠关键化合物};
\node[box,text width=44mm,h=25mm] (chemval) at (60,114) {\bt{化学合理性}%
附件8原始来源和缺陷标签；风格内关联与气味阈值\\
解释标为事后；作者加标实验作独立核验};
\node[io,text width=44mm,h=25mm] (out) at (111,114) {\bt{输出}%
49个维度各列3个候选、关联强度、区间、入选频率及证据等级\\
化合物模块供问题三使用};
\bus{boot,compare}{grade,chemval,out}{110}{32.5}{134.5}
\node[box,text width=146mm,h=13mm] (boundary) at (9,149) {\bt{结论边界}%
共表达和感官秩关联只能支持“共同变化”；第2、3位不稳定时仍按题意列出，但须标明证据不足。};
\bus{grade,chemval,out}{boundary}{145}{32.5}{134.5}
\end{tikzpicture}
\end{document}
